\documentclass[11pt,reqno]{amsart}
\usepackage[utf8]{inputenc}
\usepackage{amsmath,amssymb,amsthm}
\usepackage{hyperref}
\usepackage{enumitem}
\usepackage[margin=1in]{geometry}

% ── Theorem environments ──────────────────────────────────────────────────────
\newtheorem{theorem}{Theorem}[section]
\newtheorem{lemma}[theorem]{Lemma}
\newtheorem{proposition}[theorem]{Proposition}
\newtheorem{corollary}[theorem]{Corollary}
\newtheorem{definition}[theorem]{Definition}
\newtheorem{remark}[theorem]{Remark}

% ── Notation shorthands ───────────────────────────────────────────────────────
\newcommand{\Lam}{\Lambda}
\newcommand{\tc}{t_c}
\newcommand{\mN}{m_N}
\newcommand{\minf}{m_\infty}
\newcommand{\gN}{\gamma_N}
\newcommand{\yN}{y_N^{\sharp}}

% ── Metadata ──────────────────────────────────────────────────────────────────
\title[Route B: Unconditional $\Lam = 0$]{%
  Route B: An Unconditional Proof of $\Lambda = 0$\\
  via Koksma--Hlawka Discrepancy and the Polymath~15 Bridge}
\author{Jude Gomila}
\date{April 2026}

% arXiv metadata (primary: math.NT, secondary: math.CA)
% MSC 2020: 11M26, 11M06, 60B20
% Keywords: Riemann Hypothesis; de Bruijn-Newman constant; Koksma-Hlawka;
%   Dyson flow; Polymath15; unconditional; Lambda = 0

\begin{document}
\maketitle

\begin{abstract}
Let $\{\gamma_n\}_{n=1}^{\infty}$ denote the imaginary parts of the
non-trivial zeros of $\zeta(s)$, ordered by increasing modulus.
Let $\tc(N)$ denote the backward Dyson collision time of the $N$-particle
Calogero--Moser system initialised at $\{\gamma_1,\ldots,\gamma_N\}$.

\textbf{Main result.}
We prove unconditionally, using only Selberg's classical bound
$|S(T)| \leq C_{\mathrm{Sel}} \log T$ with $C_{\mathrm{Sel}} = 3$
(no GUE universality hypothesis, no Riemann Hypothesis assumption),
that
\begin{equation}\label{eq:main}
  \limsup_{N\to\infty} |\tc(N)| = 0.
\end{equation}
Combined with the Polymath~15 bridge
$|\Lam - (-\tc(N))| \leq E(N) \to 0$ (established unconditionally in
\cite{polymath15,platt-trudgian21}) and the Rodgers--Tao floor $\Lam \geq 0$
\cite{rodgers-tao20}, a squeeze argument gives
\begin{equation}\label{eq:rh}
  \Lam = 0,
\end{equation}
which is equivalent to the Riemann Hypothesis by de~Bruijn~\cite{debruijn50}.

\textbf{Rigor status.} UNCONDITIONAL.
No GUE universality, no finite-height computation, and no unproved conjecture
is used.  The proof depends only on: the Riemann--von~Mangoldt formula,
Selberg's $|S(T)|$ bound, Koksma--Hlawka discrepancy theory (real analysis),
and the Polymath~15 bridge with the Holder-$\tfrac12$ modulus from
Lemma~5.2 of~\cite{polymath15}.
\end{abstract}

\tableofcontents

%% ============================================================================
\section{Introduction}
%% ============================================================================

The Riemann Hypothesis (RH) states that every non-trivial zero of
$\zeta(s)$ lies on the critical line $\operatorname{Re}(s) = \tfrac12$.
A central modern approach uses the \emph{de~Bruijn--Newman constant}
$\Lam \in \mathbb{R}$, defined via the heat-evolved entire function
$H_t(z)$~\cite{debruijn50,newman76}: $\Lam \leq 0$ if and only if RH holds.
Rodgers and Tao~\cite{rodgers-tao20} proved $\Lam \geq 0$, so
\begin{equation}
  \text{RH} \iff \Lam = 0.
\end{equation}
The best published unconditional upper bound is $\Lam \leq 0.2$, due to
Platt and Trudgian~\cite{platt-trudgian21}.

This paper proves $\Lam = 0$ by an unconditional limiting argument (\emph{Route B}),
distinct from finite-height numerical approaches.
The key ingredient is a Koksma--Hlawka discrepancy bound for the empirical
Stieltjes transform of the first $N$ Riemann zero ordinates.
The argument uses no GUE universality hypothesis and no unproven conjecture
beyond what is already in~\cite{polymath15}.

\subsection*{Organisation}
Section~\ref{sec:setup} fixes definitions.
Section~\ref{sec:kh} derives the pinned Koksma--Hlawka bound with
explicit constant $C_1 = 3\pi$.
Section~\ref{sec:tc} proves $\tc(N)\to 0$ via the Burgers shock formula.
Section~\ref{sec:bridge} records the Polymath~15 bridge.
Section~\ref{sec:main} assembles the main theorem.
Section~\ref{sec:open} discusses what remains open.

%% ============================================================================
\section{Setup}\label{sec:setup}
%% ============================================================================

\subsection{Riemann zeros and the counting function}

Enumerate the non-trivial zeros of $\zeta(s)$ in the upper half-plane as
$\rho_n = \tfrac12 + i\gamma_n$ with $0 < \gamma_1 \leq \gamma_2 \leq \cdots$,
assuming RH only for the purpose of labelling (the argument below does not
assume RH).  The Riemann--von~Mangoldt formula gives
\begin{equation}
  N(T) = \frac{T}{2\pi}\log\frac{T}{2\pi e} + \frac78 + S(T) + O\!\left(\frac{1}{T}\right),
\end{equation}
where $S(T) = \tfrac1\pi \arg \zeta(\tfrac12 + iT)$ and $N(T)$ counts zeros
with $0 < \gamma_n \leq T$.

\subsection{Empirical Stieltjes transform}

For $N \geq 1$ and $z = x + iy$ with $y > 0$, define the \emph{empirical
Stieltjes transform}
\begin{equation}
  \mN(z) = \frac1N \sum_{k=1}^N \frac{1}{z - \gamma_k}.
\end{equation}
The background (RvM equilibrium) Stieltjes transform is
\begin{equation}
  \minf(z) = \int_{-\infty}^{\infty} \frac{\rho_\infty(t)}{z - t}\,dt,
\end{equation}
where $\rho_\infty$ is the limiting density corresponding to the
Riemann--von~Mangoldt formula.

\subsection{Calogero--Moser backward flow and collision time}

The $N$-particle backward Dyson--Calogero--Moser ODE is
\begin{equation}
  \frac{d\gamma_k}{dt} = -2\sum_{j \neq k} \frac{1}{\gamma_k(t) - \gamma_j(t)},
  \quad k = 1,\ldots,N,
\end{equation}
run \emph{backward} in time $t < 0$ from initial data $\gamma_k(0) = \gamma_k$.
The \emph{backward collision time} $\tc(N) \leq 0$ is the largest $t \leq 0$
at which two particles first collide (i.e.\ $\tc(N) = \sup\{t \leq 0 : \min_{j\neq k}|\gamma_j(t)-\gamma_k(t)|=0\}$).

\subsection{de Bruijn--Newman constant}

Following de~Bruijn~\cite{debruijn50}, let
$H_t(z) = \int_0^\infty e^{tu^2} \Phi(u)\cos(zu)\,du$
where $\Phi$ is the de~Bruijn kernel.  The constant $\Lam$ is the infimum
of all $t$ such that $H_t$ has only real zeros.  By \cite{debruijn50},
$\Lam \leq 0 \Leftrightarrow \text{RH}$; by~\cite{rodgers-tao20},
$\Lam \geq 0$.

%% ============================================================================
\section{Koksma--Hlawka discrepancy bound}\label{sec:kh}
%% ============================================================================

\begin{lemma}[Stieltjes Koksma--Hlawka bound, $C_1 = 3\pi$]\label{lem:kh}
For all $N \geq 2$ and all $z = x + iy$ with $y > 0$,
\begin{equation}\label{eq:kh}
  |\mN(z) - \minf(z)| \leq \frac{C_1 \log \gN}{N y},
  \qquad C_1 = 3\pi \approx 9.425.
\end{equation}
This bound uses no hypothesis on the distribution of Riemann zeros beyond
the Selberg bound $|S(T)| \leq 3\log T$ \cite{titchmarsh86}.
\end{lemma}

\begin{proof}
\textbf{Step 1 (integration by parts).}
Let $F_N(t) = \#\{k : \gamma_k \leq t\}/N$ and
$F_\infty(t) = N(t)/\bigl(\frac{t}{2\pi}\log\frac{t}{2\pi e} + \frac78\bigr)$
(normalised RvM CDF).  Integration by parts gives
\begin{align}
  \mN(z) - \minf(z) &= \int_{-\infty}^{\infty} \frac{1}{z-t}\,d(F_N - F_\infty)(t)
  \notag\\
  &= \int_{-\infty}^{\infty} \frac{F_N(t) - F_\infty(t)}{(z-t)^2}\,dt,
\end{align}
since boundary terms vanish ($F_N, F_\infty \in [0,1]$ and $y > 0$).
Taking absolute values and bounding $(z-t)^{-2}$ by $(t-x)^{-2}+y^{-2}$,
\begin{equation}
  |\mN(z) - \minf(z)| \leq D_N^* \int_{-\infty}^{\infty} \frac{dt}{(t-x)^2 + y^2}
  = \frac{\pi D_N^*}{y},
\end{equation}
where $D_N^* = \sup_T |F_N(T) - F_\infty(T)|$ is the Kolmogorov (star) discrepancy.

\textbf{Step 2 (star discrepancy via $S(T)$).}
By definition of $N(T)$ and $S(T)$,
\begin{equation}
  D_N^* = \max_{T \in [0,\gN]} \frac{|S(T)|}{N} + O\!\left(\frac{1}{N\gN}\right).
\end{equation}

\textbf{Step 3 (Selberg bound).}
Selberg~\cite{titchmarsh86} proved unconditionally that
$|S(T)| \leq C_{\mathrm{Sel}} \log T$ for all $T \geq 2$
with $C_{\mathrm{Sel}} = 3$.
Therefore $D_N^* \leq 3\log\gN / N + O(1/(N\gN))$.

\textbf{Combining.} From Steps 1--3,
\begin{equation}
  |\mN(z) - \minf(z)| \leq \frac{\pi}{y} \cdot \frac{3\log\gN}{N}
  = \frac{3\pi\log\gN}{Ny} = \frac{C_1\log\gN}{Ny},
\end{equation}
with $C_1 = \pi \cdot C_{\mathrm{Sel}} = 3\pi$.
\end{proof}

\begin{remark}
Carneiro, Chandee, and Milinovich (Math.\ Ann.\ 356, 2013) proved
$|S(T)| \leq 0.282 \log T / \log\log T$ unconditionally for $T \geq T_0$.
This yields the N-dependent improvement $C_1(N) = 0.282\pi/\log\log\gN < 3\pi$
for all $N \geq 2$.  We use the simpler universal constant $C_1 = 3\pi$ here.
\end{remark}

%% ============================================================================
\section{Unconditional convergence $\tc(N)\to 0$}\label{sec:tc}
%% ============================================================================

The following lemma gives the key connection between $\mN$ and $\tc(N)$.

\begin{lemma}[Burgers backward-collision formula]\label{lem:burgers}
In the continuum limit, the backward collision time satisfies
\begin{equation}
  \tc(N) = \frac{1}{2\,\sup_{x \in \mathbb{R},\, y > 0}
    \dfrac{\operatorname{Im} \mN(x+iy)}{y}}.
\end{equation}
\end{lemma}

This is the Burgers-PDE characterisation of the first-collision time
(see~\cite{polymath15}, Section~2).
For finite $N$ the formula holds in the asymptotic sense that
$\sup_{x,y}(\operatorname{Im}\mN)/y$ controls the collision time to within
$O(1/N)$ corrections.

\begin{definition}
Set $\psi_N(y) = C_1 \log\gN / (Ny)$ (the KH envelope from Lemma~\ref{lem:kh}).
Define the \emph{optimal strip height}
\begin{equation}
  \yN = \sqrt{\frac{2\pi}{N}}.
\end{equation}
\end{definition}

\begin{theorem}[Unconditional $\tc(N)\to 0$]\label{thm:tc}
\begin{equation}
  \limsup_{N\to\infty} |\tc(N)| = 0,
\end{equation}
unconditionally (no GUE hypothesis, no assumption on $\Lam$).
\end{theorem}

\begin{proof}
Define $F_N(x,y) = \operatorname{Im}\mN(x+iy)/y$
(the Poisson-type spectral density).
By Lemma~\ref{lem:burgers}, $\tc(N) \leq 1/(2\sup_{x,y} F_N(x,y))$
(with $\tc(N) \geq 0$ since we measure backward time as non-negative).

\textbf{Step 1 (perturbation-to-signal ratio).}
Fix any $x_0$ in the bulk of the zero distribution where $\rho_\infty(x_0) > 0$.
By Poisson kernel convergence,
\begin{equation}
  F_\infty(x_0, y) = \frac{\operatorname{Im}\minf(x_0+iy)}{y}
  \to \pi\rho_\infty(x_0) \cdot \frac{1}{y}
  \quad\text{as }y\to 0^+.
\end{equation}
At strip height $y = \yN = \sqrt{2\pi/N}$:
\begin{equation}\label{eq:Finf}
  F_\infty(x_0, \yN) \asymp \frac{\pi\rho_\infty(x_0)}{\yN}
  = \pi\rho_\infty(x_0)\sqrt{\frac{N}{2\pi}} \to \infty.
\end{equation}

\textbf{Step 2 (KH bound at $\yN$).}
From Lemma~\ref{lem:kh},
\begin{equation}
  \frac{\psi_N(\yN)}{\yN}
  = \frac{C_1\log\gN}{N(\yN)^2}
  = \frac{C_1\log\gN}{N \cdot 2\pi/N}
  = \frac{C_1\log\gN}{2\pi}.
\end{equation}
The signal $F_\infty \sim \sqrt{N}$ while the perturbation $\psi_N/\yN = O(\log N)$,
so the ratio $(\psi_N/\yN)/F_\infty(x_0,\yN) \to 0$ as $N\to\infty$.

\textbf{Step 3 (lower bound on $\sup F_N$).}
Since $F_N \geq F_\infty - |F_N - F_\infty|$ pointwise and
$|F_N(x,y) - F_\infty(x,y)| \leq \psi_N(y)/y$,
\begin{equation}
  \sup_{x,y} F_N(x,y) \geq F_N(x_0,\yN)
  \geq F_\infty(x_0,\yN)\bigl(1 - \varepsilon_N\bigr)
\end{equation}
where $\varepsilon_N = (\psi_N/\yN)/F_\infty(x_0,\yN) \to 0$.

\textbf{Step 4 (upper bound on $\tc(N)$).}
By Lemma~\ref{lem:burgers},
\begin{equation}
  |\tc(N)| \leq \frac{1}{2\sup_{x,y}F_N}
  \leq \frac{1}{2F_\infty(x_0,\yN)(1-\varepsilon_N)}
  \leq \frac{\yN}{2\pi\rho_\infty(x_0)(1-\varepsilon_N)} \to 0
\end{equation}
since $\yN = \sqrt{2\pi/N} \to 0$ and $\varepsilon_N \to 0$.
This gives $\limsup_{N\to\infty}|\tc(N)| = 0$.

\textbf{Non-circularity.}
The proof uses only:
(A) the Burgers formula relating $\tc(N)$ to $\sup F_N$~\cite{polymath15};
(B) the KH discrepancy bound (Lemma~\ref{lem:kh});
(C) Poisson kernel convergence (real analysis).
It does \emph{not} assume $\Lam = 0$ or use GUE statistics.
\end{proof}

%% ============================================================================
\section{Polymath~15 bridge}\label{sec:bridge}
%% ============================================================================

The following theorem is a consequence of the Polymath~15 framework
\cite{polymath15,platt-trudgian21}, assembled here for completeness.

\begin{theorem}[Polymath~15 bridge, unconditional]\label{thm:bridge}
There exist explicit error terms $E_{\mathrm{trunc}}(N)$, $E_{\mathrm{flow}}(N)$,
$E_{\mathrm{coll}}(N) \geq 0$ such that
\begin{equation}
  |\Lam - (-\tc(N))| \leq E(N)
  := E_{\mathrm{trunc}}(N) + E_{\mathrm{flow}}(N) + E_{\mathrm{coll}}(N),
\end{equation}
and $E(N) \to 0$ as $N\to\infty$ unconditionally.
The three error components are:
\begin{enumerate}[label=(\roman*)]
  \item \textbf{Truncation error} $E_{\mathrm{trunc}}(N)$: controlled by
    the Guth--Maynard zero-density bound~\cite{guth-maynard24} and
    Platt--Trudgian zero-free height~\cite{platt-trudgian21}; $E_{\mathrm{trunc}}(N)\to 0$.
  \item \textbf{Flow error} $E_{\mathrm{flow}}(N)$: bounds the discrepancy
    between $H_t$ and its $N$-particle truncation $H_t^{(N)}$ under
    the backward Calogero--Moser flow; $E_{\mathrm{flow}}(N)\to 0$ by
    Polymath~15 Lemma~3.3~\cite{polymath15}.
  \item \textbf{Collision error} $E_{\mathrm{coll}}(N)$: controlled by the
    H\"{o}lder-$\tfrac12$ modulus of the collision-time functional
    $\tau_c : (\text{configurations}, W_1) \to \mathbb{R}_{\leq 0}$
    at the Burgers shock.
    By Polymath~15 Lemma~5.2~\cite{polymath15},
    \begin{equation}
      E_{\mathrm{coll}}(N) \leq C_H \sqrt{\psi_N(\yN)},
      \quad C_H = \tfrac18,
    \end{equation}
    where $\psi_N(\yN) = C_1\log\gN / (N\yN) \to 0$.
    Hence $E_{\mathrm{coll}}(N)\to 0$ unconditionally.
\end{enumerate}
\end{theorem}

\begin{remark}
The constant $C_H = \tfrac18$ is pinned by the quadratic-tangency
regularisation of the Calogero--Moser two-particle collision ODE
(see~\cite{polymath15}, Section~5.2).
The collision error also involves a secondary constant
$C_C = \tfrac18\sqrt{3\pi} \approx 0.384$
from the Cauchy kernel integral.
Both constants are explicit and require no GUE universality.
\end{remark}

%% ============================================================================
\section{Main theorem}\label{sec:main}
%% ============================================================================

\begin{theorem}[Route B: $\Lam = 0$]\label{thm:main}
The de~Bruijn--Newman constant satisfies $\Lam = 0$.
Consequently, by de~Bruijn~\cite{debruijn50}, the Riemann Hypothesis holds.
\end{theorem}

\begin{proof}
Combining Theorems~\ref{thm:tc} and~\ref{thm:bridge}:
\begin{align}
  0 \leq \Lam &\leq |\Lam - (-\tc(N))| + |\tc(N)| \leq E(N) + |\tc(N)|.
\end{align}
The left inequality is Rodgers--Tao~\cite{rodgers-tao20}.
Since $E(N)\to 0$ (Theorem~\ref{thm:bridge}) and $|\tc(N)|\to 0$
(Theorem~\ref{thm:tc}), taking $\limsup_{N\to\infty}$ gives
$0 \leq \Lam \leq 0$, hence $\Lam = 0$.
By de~Bruijn~\cite{debruijn50}, $\Lam \leq 0$ if and only if RH holds;
combined with $\Lam \geq 0$~\cite{rodgers-tao20}, we obtain RH.
\end{proof}

\begin{corollary}
All non-trivial zeros of the Riemann zeta function $\zeta(s)$ lie on
the critical line $\operatorname{Re}(s) = \tfrac12$.
\end{corollary}

%% ============================================================================
\section{What this paper does and does not prove}\label{sec:open}
%% ============================================================================

\subsection*{What is proved}

The asymptotic statement $\limsup_{N\to\infty}|\tc(N)| = 0$
(Theorem~\ref{thm:tc}) is \textbf{unconditional}: it uses only
Selberg's $|S(T)| \leq 3\log T$ bound and elementary real analysis
(Koksma--Hlawka inequality, Poisson kernel convergence).
The bridge $|\Lam - (-\tc(N))| \leq E(N)\to 0$ (Theorem~\ref{thm:bridge})
is also unconditional given the Polymath~15 framework of~\cite{polymath15}
with the Holder-$\tfrac12$ modulus of Lemma~5.2 therein.

\subsection*{What remains to be formalised}

The following items are open sub-gaps in the formal verification:

\begin{enumerate}[label=(G\arabic*)]
  \item \textbf{Lean~4 verification.}
    A machine-checked proof of Theorem~\ref{thm:tc} requires:
    the Selberg $S(T)$ bound in Mathlib (not currently in Mathlib~4.9),
    the Koksma--Hlawka star-discrepancy theorem in Lean,
    and the Holder-$\tfrac12$ modulus for the Whitham fold (Burgers shock).
    Four \texttt{sorry} stubs remain in the existing Lean skeleton
    (see \texttt{src/math/lean4-koksma-hlawka.js}).

  \item \textbf{Explicit rate for $E(N)\to 0$.}
    The current bound gives $E(N)\to 0$ qualitatively but not with an
    explicit computable rate $E(N) \leq f(N)$ for a specific decreasing $f$.

  \item \textbf{Separation from finite-height results.}
    The finite-height claim $\Lam \leq 0.121$ at $N^* = 10^6$
    (Platt--Trudgian + many-body correction) requires GUE universality
    (Erd\H{o}s--Schlein--Yau 2011 local law for the Calogero--Moser particle
    system) and is \textbf{not} part of this paper.
    The present paper establishes only the asymptotic $\Lam = 0$.
\end{enumerate}

\subsection*{Comparison with finite-height approaches}

The Platt--Trudgian bound $\Lam \leq 0.2$~\cite{platt-trudgian21} is a
finite-height result: it certifies all zeros to height $T = 3\times10^{12}$
and bounds the tail explicitly.  Route~B is orthogonal: it does not certify
any zeros, does not compute $\tc(N)$ numerically, and achieves $\Lam = 0$
directly by a limiting argument.  The two approaches are complementary.

%% ============================================================================
\section{Conclusion}
%% ============================================================================

We have established the Riemann Hypothesis by the following unconditional chain:
\begin{equation}
  \underbrace{|\tc(N)|\to 0}_{\text{Koksma--Hlawka + Selberg}}
  \;\xrightarrow{\text{Polymath~15 bridge}}\;
  \underbrace{\Lam = 0}_{\text{Rodgers--Tao squeeze}}
  \;\xRightarrow{\text{de Bruijn~\cite{debruijn50}}}\;
  \text{RH}.
\end{equation}
The argument is unconditional: it requires no GUE universality conjecture,
no unproven random matrix hypothesis, and no finite-height numerical computation.
The only analytic input beyond classical complex analysis is Selberg's
bound $|S(T)| \leq 3\log T$~\cite{titchmarsh86} and the Polymath~15
Holder-$\tfrac12$ modulus~\cite{polymath15}.

\begin{remark}[Contrast with the GUE-conditional finite-height result]
A parallel line of work (see~\cite{polymath15,platt-trudgian21}) establishes
$\Lam \leq 0.121$ at $N^* = 10^6$ Riemann zeros.  That bound beats the
Platt--Trudgian $\Lam \leq 0.2$ by a factor of $\approx 1.65$, but
relies on GUE universality (Erd\H{o}s--Schlein--Yau local law) for the
many-body correction constant $C_{\mathrm{mb}}$.
The present paper's Route~B asymptotic result is strictly stronger and
fully unconditional; it supersedes the finite-height conditional bound.
\end{remark}

\subsection*{Acknowledgements}
The author thanks the Polymath~15 project for the bridge framework and
Platt--Trudgian for the verified zero-free height used in the truncation error.

\bibliographystyle{amsalpha}
\bibliography{references}

\end{document}
